\documentclass[12pt,a4paper]{article}
\usepackage[latin1]{inputenc}
\usepackage{amsmath}
\usepackage{amsfonts}
\usepackage{amssymb}
\usepackage{makeidx}
\usepackage{graphicx,textcomp,marvosym}
\usepackage[left=3.00cm, right=3.00cm, top=4.00cm, bottom=3.00cm]{geometry}
\author{%Parley Ruogu Yang\footnote{Department of Statistics, University of Oxford and Faculty of Mathematics, University of Cambridge}, Hugo Dolan\footnote{University College Dublin}, Ondrej Bohdal\footnote{School of Informatics, University of Edinburgh} \\ 
Team HOP for Luno \& General Challenge, Spark Hackathon }
\title{Trading Strategies Fact Sheet}
\date{This version: 14 Aug 2020}
\begin{document}
\thispagestyle{empty}
\maketitle
\vfill
\begin{figure}[h]
	\centering
	\includegraphics[width=\linewidth]{../Logo}
\end{figure}
\vfill
\tableofcontents
\thispagestyle{empty}
	
	\pagebreak
	
\vfill
\section{Trading Summary}
\begin{itemize}
\item We designed two Market Making Strategies (M1 \& M2) which can facilitate potential liquidity provision in Luno's exchange and make low-risk profits.
\item We designed two Speculative Strategies (S1 \& S2) which apply traditional trading strategies that can  make profits with controllable risks.
\item We translated M1, M2, and S1 into the LIMA system and deployed them successfully in a live trading environment.
\item However we soon discovered that the bid-ask spread of certain currency pairs (e.g. BTC/EUR) on many exchanges including LUNO is quite tight meaning that the margin on each trade is quite small. Compounding this effect we found that trades occurring on LUNO tend to be quite small in volume and also relatively infrequent compared to larger exchanges. This particularly affected our statistical arbitrage strategies (M1 and in some cases  S1) as they relied on using limit orders with lower execution costs to achieve a profit, thus low levels of liquidity on LUNO made these strategies infeasible as many arbitrage opportunities exist for shorter periods than the average time to fill on LUNO.

\item Similarly we attempted a liquidity provision strategy (M2) which increased the trade volume on LUNO thus aimed to capture a small profit from executing the market and limit orders on both sides of the book. This strategy was also successfully executed, however due to the costs of LUNO market orders the profitability of these strategies is unclear as of yet,  as we have only run this particular strategy for several days.

\item From our trading work, we suggest that further research should focus more on the longer term (hourly or daily) technical trading strategies rather than shorter term market making style approach taken in our existing research. This would provide the opportunity for larger profits per trade due to possible large intra-day market volatility and eliminate concerns about trading costs.

\end{itemize}

\section{Introductory Background}
 \begin{itemize}
		\item There are some degrees of lead-lag relationships between the main exchange (i.e. the more popular ones, here we use Bitstamp and Kraken for example) and the Luno's prices, especially on crypto-sovereign currency pairs. This is mainly due to the lack of liquidity on Luno's exchange.
		\item There are arbitrage opportunities both in a crypto-sovereign-crypto currency route and a crypto-sovereign-sovereign route. 
		\item Traditional financial strategies are functional, subject to trading costs, which are zero for limit orders but 0.25\% for filling the limit orders or market orders.
	
	\end{itemize}
\vfill
	\pagebreak
	\section{Technical details of the strategies}
		\subsection{Market Making Strategy 1 (M1)}
		\begin{table}[h]
			\centering
		\begin{tabular}{c|c}
			
		Profitability& \textdollaroldstyle \\
		
		Risk & \Radioactivity \\
		
		Liquidity provision &$ \bigtriangleup \bigtriangleup $\\
		
		\end{tabular}
		\end{table}
	
This is a classic market-making / statistical arbitrage strategy where we buy from low and sell from high the difference on the same asset.

At time $t$, write $y_t$ as the LUNO price of an asset and $z_t$ as the main exchange price.
	Let $C_1$ be the threshold of equal, whereby we judge $|y_t-z_t|<C_1$ as ``equal state".
\begin{itemize}
	\item 	If we start from being equal to a state where $y_t>z_t+C_1$, then we initiate a trade by shorting LUNO and longing from the main exchange and aim at a price level near $z_t+C_1$ as we believe $y_t$ will fall to such a level soon.
	\item If we start from being equal to a state where  $y_t< z_t -C_1$, then we initiate a trade by longing LUNO and shorting from the main exchange and aim at a price level near $z_t-C_1$ as we believe $y_t$ will raise to such a level soon.
\end{itemize}

\textbf{Note on execution}: We execute them in a Fill-or-Kill limit order manner on LUNO, i.e. we provide one-sided liquidity to LUNO and only trigger the strategy if our limit order gets filled. The execution on main exchange is of no concern, benefited from its liquidity.
	
\textbf{Note on LUNO's asset}: The simple asset can be BTC/GBP and BTC/EUR while we may extend this definition to a chained asset. For example, while $y_t$ being BTC/GBP, we may have $z_t$ being BTC/EUR $\times$ EUR/GBP.


\pagebreak
\subsection{Market Making Strategy 2 (M2)}

		\begin{table}[h]
	\centering
	\begin{tabular}{c|c}
		
		Profitability& \textdollaroldstyle \\
		
		Risk & \Radioactivity  \\
		
		Liquidity provision &$ \bigtriangleup \bigtriangleup  \bigtriangleup$\\
		
	\end{tabular}
\end{table}

This is a Limit-Order-Book (LOB) based market-making strategy in which we aim to provide liquidity on both side of the order book and hedge our position from the main exchange.

At time $t$, write $P^{bid}_t$ and $P^{ask}_t$ as the best bid and ask prices for an asset on LUNO's exchange (i.e. LUNO's LOB). Write $z_t$ as the mid price of the same asset on the main exchange.

Consider the spread $S_t := P^{ask}_t - P^{bid}_t$. We aim to provide liquidity by shrinking the spread by $S_t^\prime >0$, thus start by putting a limit sell order at $P^{ask}_t - \frac{S_t^\prime}{2}$ and a limit buy order at $P^{bid}_t + \frac{S_t^\prime}{2}$.

Upon triggering any of the order, we enter a risk management --- exit state, where: \begin{itemize}
	\item If $z_t$ moves in the adverse direction,\footnote{E.g. if we are long the asset and the price moves down.} we stop loss at an appropriate level and execute at the main exchange to secure the price. Full exit happens when we close our current position on LUNO via limit order and the hedged position on the main exchange.
	\item If $z_t$ moves in the favourable direction, \footnote{E.g. if we are long the asset and the price moves up.} we keep offering the best side of the book to be filled, that is, if we are long we keep posting a limit sell order at $P^{ask}_t - \varepsilon$ with a small $\varepsilon > 0$, and vice versa for the case if we are short.
\end{itemize}



\vfill

\pagebreak

\subsection{Speculative Strategy 1 (S1)}

\begin{table}[h]
	\centering
	\begin{tabular}{c|c}
		
		Profitability& \textdollaroldstyle \textdollaroldstyle \textdollaroldstyle \\
		
		Risk & \Radioactivity  \Radioactivity \Radioactivity \\
		
		Liquidity provision &$ \bigtriangleup$\\
		
	\end{tabular}
\end{table}

In this strategy we chase the break-out / momentum via traditional statistical indicators on the main exchange.

At time $t$, write $z_t$ as the mid price of the same asset on the main exchange. We wait for a trigger, either up- or down-side of break-out on $z_t$, this can be, for example, a penetration of a Bollinger Band or other time series evaluators.

We buy or sell via market order first to obtain the entry\footnote{If the LUNO's spread is too wide, this would have to be done in the main exchange with the possibility to transfer it back to LUNO at an appropriate time.}, then place limit order at the associate target price. E.g. if we use an upper $2\sigma$ Bollinger as the break-out signal, we buy via market order as the price penetrates the upper Bollinger, then wait until, say the  $3\sigma$ Bollinger or higher for an exit via limit order.

Risk management is done by setting an associated exit handle, whereas if $z_t$ drops below such level, a stop-market order to close the position is entered at the main exchange with a transfer back to LUNO at an appropriate time.
%\vfill
%\pagebreak
\subsection{Speculative Strategy 2 (S2)}

\begin{table}[h]
	\centering
	\begin{tabular}{c|c}
		
		Profitability& \textdollaroldstyle  \textdollaroldstyle \\
		
		Risk & \Radioactivity  \Radioactivity \\
		
		Liquidity provision &$ \bigtriangleup \bigtriangleup$\\
		
	\end{tabular}
\end{table}
This is a more rewarding --- thus more risky version of M2, where we do not execute on the main exchange.

We inherit the same notation as in M2 and use $z_t$ as a signal feed-in. But, when we execute the exit, we always do it in a limit sell order at $P^{ask}_t - \varepsilon$. This decreases the trading complexity thus cost, but is more risky in chaotic market conditions.

An additional option is we exit by stages --- this reduces the risk while ensures the liquidity provision on one side to be longer. Another variation is that we do multiple stages of entry --- this can be achieved by having, say, multiple 
pairs of $\{(P^{ask}_t - \frac{S_t^\prime}{2} , P^{bid}_t + \frac{S_t^\prime}{2}) | S_t^\prime \in X \}$ for a set of spread improvements $X$ and wait for execution and its associated staggered exit strategies. 



%\pagebreak
%\section{Key Findings from backtestting \& trading}


%Lack of liquidity on LUNO: the rate of limit orders being filled on LUNO is slow, thus profitability in some strategies are lowered as we may be forced to execute market orders.

\end{document}